% ==================================================================
% id:       prob-015
% topic:    מבוא להסתברות
% subtopic: מרחב הסתברות אחיד
% solution: full
% ==================================================================
בכד יש 7 כדורים שחורים וכדור צהוב אחד. בוחרים מהכד באקראי 3 כדורים (בבת אחת). מה ההסתברות שהכדור הצהוב נבחר? (אפשר להבחין בין כדורים מאותו צבע.)
\begin{solbox}
נמספר את הכדורים. מרחב המדגם הוא כל תתי-הקבוצות בגודל 3 של 8 הכדורים, כולן סבירות באותה מידה, ו-$|\Omega|=\binom83=56$.

\textbf{דרך א.} בתוצאות שבהן הצהוב נבחר צריך לבחור עוד 2 כדורים מתוך 7 השחורים: $\binom72=21$. לכן ההסתברות היא $\frac{21}{56}=\ans{\frac38}$.

\textbf{דרך ב: המשלים.} בתוצאות שבהן הצהוב \emph{לא} נבחר, שלושת הכדורים שחורים: $\binom73=35$. לכן ההסתברות היא $1-\frac{35}{56}=\frac38$.
\end{solbox}
